\documentclass[a4paper]{article}
\usepackage{amsmath, amsthm, amssymb, bm, graphicx, hyperref, mathrsfs}
\usepackage{booktabs}
\usepackage[affil-it]{authblk}
\usepackage[backend=bibtex,style=numeric]{biblatex}

\usepackage{geometry}
\geometry{margin=1.5cm, vmargin={0pt,1cm}}
\setlength{\topmargin}{-1cm}
\setlength{\paperheight}{29.7cm}
\setlength{\textheight}{25.3cm}

\addbibresource{citation.bib}

\begin{document}
% =================================================
\title{Numerical Analysis Homework XXX}

\author{Yilin Chen 3230101092
  \thanks{Electronic address: \texttt{yilinchen@zju.edu.cn}}}
\affil{(Numerical Analysis, Math2301), Zhejiang University }


\date{Due time:25/11/21 \today}

\maketitle

% \begin{abstract}
%     The abstract is not necessary for the theoretical homework, 
%     but for the programming project, 
%     you are encouraged to write one.      
% \end{abstract}





% ============================================
\section*{I. Convert a decimal integer to a normalzed FPN}

Using Algorithm 4.8, calculate 

$477 / 2 = 238 ... 1$

$238 / 2 = 119 ... 0$

$119 / 2 = 59 ... 1$

$59 / 2 = 29 ... 1$

$29 / 2 = 14 ... 1$

$14 / 2 = 7 ... 0$

$7 / 2 = 3 ... 1$

$3 / 2 = 1 ... 1$

$1 / 2 = 0 ... 1$

we know $477 = (111011101)_2 = (1.11011101) \times 2^{8}$


\section*{II. Convert a decimal fraction to a normalzed FPN}

Using Algorithm 4.8, calculate

$0.6 \times 2 = 1.2$, where the integer is $1$ and the decimal is $0.2$ 

$0.2 \times 2 = 0.4$, where the integer is $0$ and the decimal is $0.4$

$0.4 \times 2 = 0.8$, where the integer is $0$ and the decimal is $0.8$

$0.8 \times 2 = 1.6$, where the integer is $1$ and the decimal is $0.6$

$0.6 \times 2 = 1.2$, where the integer is $1$ and the decimal is $0.2$ (repeat)

Thus, $\dfrac{3}{5} = (0.10011001\cdots)_2 = (1.00110011\cdots)_2 \times 2^{-1}$

\section*{III. Adjacent FPNs}

Since \( x \) is a normalized FPN, it can be written as:
\[
x = m \times \beta ^e, \quad \text{where } m = 1.
\]
The set of normalized FPNs in \( \mathcal F \) with exponent \( e \) consists of numbers of the form:
\[
m \times \beta^e, \quad \text{where } m = \frac{n}{\beta^{p-1}}, \quad n \in \{\beta^{p-1}, \beta^{p-1}+1, \dots, \beta^p - 1\}.
\]
The spacing between consecutive numbers in this exponent range is:
\[
\Delta_e = \frac{1}{\beta^{p-1}} \cdot \beta^e = \beta^{e-p+1}.
\]
Thus, the next number after \( x \) in exponent \( e \) is:
\[
Z_R = \left(1 + \frac{1}{\beta^{p-1}}\right) \beta^e = \beta^e + \beta^{e-p+1}.
\]

Now, consider numbers with exponent \( e-1 \). The normalized FPNs with exponent \( e-1 \) are:
\[
m \times \beta^{e-1}, \quad \text{with } m = \frac{n}{\beta^{p-1}}, \quad n \in \{\beta^{p-1}, \beta^{p-1}+1, \dots, \beta^p - 1\}.
\]
The largest number in this set is:
\[
M_{\text{max}} \times \beta^{e-1} = \frac{\beta^p - 1}{\beta^{p-1}} \cdot \beta^{e-1} = \left(\beta - \frac{1}{\beta^{p-1}}\right) \beta^{e-1} = \beta^e - \beta^{e-p}.
\]
This number is less than \( x \), and there is no normalized FPN between this number and \( x \) because any number with exponent \( e-1 \) is at most this value, and any number with exponent \( e \) is at least \( x \). Therefore,
\[
Z_L = \beta^e - \beta^{e-p}.
\]

Then we compute:
\[
x - Z_L = \beta^e - (\beta^e - \beta^{e-p}) = \beta^{e-p},
\]
and
\[
Z_R - x = (\beta^e + \beta^{e-p+1}) - \beta^e = \beta^{e-p+1}.
\]
Thus,
\[
Z_R - x = \beta^{e-p+1} = \beta \cdot \beta^{e-p} = \beta (x - Z_L).
\]
This completes the proof.


\section*{IV. Adjacent FPNs for \( x = 3/5 \) under IEEE 754 Single-Precision}

Under the IEEE 754 single-precision protocol, the parameters are:
\begin{itemize}
\item Base \(\beta = 2\)
\item Precision \(p = 24\) (including the implicit leading bit)
\item Exponent range: \(L = -126\), \(U = 127\)
\end{itemize}

From II, we know \(\dfrac35 = (1.001100110011\cdots)_2 \times 2^{-1}\). Thus, the exponent is \(e = -1\), and the exponent field is \(-1 + 127 = 126\), represented in binary as \(01111110\).

The significand \(s = 1.2\) in decimal, so \(s - 1 = 0.2\). The integer value for the fraction part is computed as \(k = \lfloor (s-1) \times 2^{23} \rfloor = \lfloor 0.2 \times 8388608 \rfloor = \lfloor 1677721.6 \rfloor = 1677721\). Thus, the two adjacent normalized FPNs are:
\begin{align*}
Z_L &= \left(1 + \frac{1677721}{2^{23}}\right) \times 2^{-1} \\
Z_R &= \left(1 + \frac{1677722}{2^{23}}\right) \times 2^{-1}
\end{align*}

In IEEE 754 binary format:
\begin{itemize}
\item \(Z_L\): sign bit \(0\), exponent \(01111110\), fraction \(00110011001100110011001\)
\item \(Z_R\): sign bit \(0\), exponent \(01111110\), fraction \(00110011001100110011010\)
\end{itemize}

The value \(\text{fl}(x)\) is the rounded value of \(x\) under the round-to-nearest mode. Since \(x\) is closer to \(Z_R\) than to \(Z_L\), we have \(\text{fl}(x) = Z_R\).

The absolute error is \(|Z_R - x| = 0.4 \times 2^{-24}\). The relative roundoff error is:
\[
\frac{|Z_R - x|}{|x|} = \frac{0.4 \times 2^{-24}}{0.6} = \frac{2}{3} \times 2^{-24} = \frac{2^{-23}}{3}
\]

\section*{V. }


\section*{VI. }

\section*{VII. }


\section*{VIII. }

% ===============================================
% \section*{ \center{\normalsize {Acknowledgement}} }
% Give your acknowledgements here(if any).


% \printbibliography

% If you are not familiar with \texttt{bibtex}, 
% it is acceptable to put a table here for your references.
\end{document}